\origpage{336}{N. Robertson, P.D. Seymour / Journal of Combinatorial Theory, Series B 90 (2004) 325--385}
\hspace*{12pt}$\bullet$ if $\psi(R) \subseteq S$ for some side $S$ of $\Phi$, then $R$ is directed such that $\psi$ maps its direction onto that of $S$, and $R$ is designated as long if and only if $S$ is long,

\hspace*{12pt}$\bullet$ if there is no such $S$, then $R$ is directed arbitrarily and is designated as short.

Now let $\Gamma$ be a painting fitting $\Phi$, such that $U(\Gamma) \cap U(K) = V(\Gamma) \cap V(K)$. We denote by $\psi^{-1}(\Gamma) \cap \Sigma_1$ the painting $\Gamma_1 = (U_1, V_1, \gamma_1)$ in $\Sigma_1$, where
\begin{align*}
U_1 &= (\psi^{-1}(U(\Gamma)) \cap \Sigma_1) \cup bd(\Sigma_1), \\
V_1 &= \psi^{-1}(V(\Gamma) \cup V(K)) \cap \Sigma_1
\end{align*}
and for each edge $e$ of $\Gamma_1$, if $e \subseteq bd(\Sigma_1)$ then $\gamma_1(e)$ is defined such that the direction of $e$ agrees with the direction of the side of $\psi^{-1}(\Phi) \cap \Sigma_1$ including it, while if $e \nsubseteq bd(\Sigma_1)$ then $\gamma_1(e)$ is defined such that $\psi$ maps it to $\gamma_\Gamma(\psi(e))$.

\textbf{5.3.} \textit{With $\Gamma, K$, etc. as above, if each component $K'$ of $K$ satisfies}
\begin{equation*}
|V(K') \cap V(\Gamma)| + |V(K') \cap bd(\Sigma)| \ge 2
\end{equation*}
\textit{then $\psi^{-1}(\Gamma) \cap \Sigma_1$ fits the frame $\psi^{-1}(\Phi) \cap \Sigma_1$.}

\textbf{Proof.} We observe first that the hypothesis implies that for every component $K'$ of $K$, $V(K')$ contains a vertex of $\Gamma$, since $V(K) \cap bd(\Sigma) \subseteq V(K) \cap U(\Gamma) \subseteq V(\Gamma)$. Verifying the first three conditions in the definition of ``fit'' is easy\del{ and is omitted}. \fix{We give the details, and we check also that $\Phi_1 = \psi^{-1}(\Phi) \cap \Sigma_1$ is a frame, that is, that no two long sides of $\Phi_1$ have a common end. Now $\psi$ maps $bd(\Sigma_1)$ into $bd(\Sigma) \cup U(K)$, and $bd(\Sigma) \cap U(K) \subseteq V(K)$, since $bd(\Sigma) \subseteq U(\Gamma)$ and $U(\Gamma) \cap U(K) = V(\Gamma) \cap V(K)$. Since each side $R$ of $\Phi_1$ is arcwise connected and disjoint from $\psi^{-1}(V(\Phi) \cup V(K))$, it follows that $\psi(R)$ is included either in a side of $\Phi$ or in an edge of $K$.

Suppose that two long sides $R_1, R_2$ of $\Phi_1$ have a common end $p$. Near $p$, $bd(\Sigma_1)$ consists of $p$ and points of $R_1 \cup R_2$, and so it is mapped by $\psi$ into $bd(\Sigma)$. Consequently $\psi$ maps the points of $\Sigma_1 \setminus bd(\Sigma_1)$ near $p$ onto all the points of $\Sigma \setminus bd(\Sigma)$ near $\psi(p)$; and since $\psi^{-1}(U(K)) \subseteq bd(\Sigma_1)$, no point of $U(K) \setminus bd(\Sigma)$ is near $\psi(p)$. Since $K$ has no isolated vertices and no border edges, it follows that $\psi(p) \notin V(K)$, and so $\psi(p) \in V(\Phi)$. For $i = 1, 2$, since $R_i$ is long, $\psi(R_i)$ is included in a long side $S_i$ of $\Phi$, and $\psi(p)$ is an end of $S_i$. Since no two long sides of $\Phi$ have a common end, $S_1 = S_2$. Now $\psi$ maps the points of $R_1$ and of $R_2$ near $p$ onto the points of $S_1$ near its end $\psi(p)$, and so there exist $x_1 \in R_1$ and $x_2 \in R_2$, both near $p$, with $\psi(x_1) = \psi(x_2)$. Near $x_1$ and near $x_2$, as near $p$, $\psi$ maps the points of $\Sigma_1 \setminus bd(\Sigma_1)$ onto all the points of $\Sigma \setminus bd(\Sigma)$ near $\psi(x_1)$; since $\psi$ is one-to-one on $\Sigma_1 \setminus bd(\Sigma_1)$, it follows that $x_1 = x_2$, and so $R_1 \cap R_2 \ne \emptyset$, which is impossible. Thus $\Phi_1$ is a frame.

For the first condition, $V(\Phi_1) = \psi^{-1}(V(\Phi) \cup V(K)) \cap \Sigma_1 \subseteq \psi^{-1}(V(\Gamma) \cup V(K)) \cap \Sigma_1 = V(\Gamma_1)$, since $V(\Phi) \subseteq V(\Gamma)$. For the second, let $R$ be a short side of $\Phi_1$. If $\psi(R)$ is included in a side $S$ of $\Phi$, then $S$ is short, by the definition of $\Phi_1$; so $S$ is an edge of $\Gamma$, since $\Gamma$ fits $\Phi$, and therefore $S \cap V(\Gamma) = \emptyset$. Otherwise $\psi(R)$ is included in an edge $f$ of $K$, and $f \cap V(\Gamma) = \emptyset$ since $U(\Gamma) \cap U(K) = V(\Gamma) \cap V(K)$. In both cases $R$ is disjoint from $\psi^{-1}(V(\Gamma) \cup V(K))$, and so $R$ is an arcwise connected component of $bd(\Sigma_1) \setminus V(\Gamma_1)$; that is, $R$ is an edge of $\Gamma_1$. For the third, let $e$ be a border edge of $\Gamma_1$. Every edge of $\Gamma_1$ not included in $bd(\Sigma_1)$ is of the form $\psi^{-1}(f) \cap \Sigma_1$ for some internal edge $f$ of $\Gamma$, and is disjoint from $bd(\Sigma_1)$, since $\psi$ maps $bd(\Sigma_1)$ into $bd(\Sigma) \cup U(K)$, which $f$ does not meet. Hence $e \subseteq bd(\Sigma_1)$, and $\gamma_1(e)$ was defined so that the direction of $e$ agrees with the direction of the side of $\Phi_1$ including it.}\fixnote{fix:5.3-proof-added}{The paper omits the verification of the first three conditions of ``fit'', and of the condition in the definition of a frame that no two long sides have a common end. The argument above is the one used in the formalization.} Let us verify the fourth and fifth conditions. Let $\Gamma_1 = \psi^{-1}(\Gamma) \cap \Sigma_1$, and let $F_1$ be a $\Gamma_1$-normal $O$-arc in $\Sigma_1$ with $|F_1 \cap V(\Gamma_1)| + |F_1 \cap bd(\Sigma_1)| \le 2$. Then $|F_1 \cap bd(\Sigma_1)| \le 1$ since $F_1 \cap bd(\Sigma_1) \subseteq F_1 \cap V(\Gamma_1)$, and so $F = \psi(F_1)$ is a $\Gamma$-normal $O$-arc in $\Sigma$ with $|F \cap V(\Gamma)| + |F \cap bd(\Sigma)| \le 2$. Since $\Gamma$ is internally 3-connected, there is a disc $\Delta \subseteq \Sigma$ bounded by $F$ with $\Delta \cap V(\Gamma) \subseteq F$. If $U(K) \cap (\Delta \setminus F) = \emptyset$, then $\psi^{-1}(\Delta)$ includes a disc in $\Sigma_1$ bounded by $F_1$ as required. We assume then (for a contradiction) that $U(K) \cap (\Delta \setminus F) \ne \emptyset$. Let $K'$ be a component of $K$ with $U(K') \cap (\Delta \setminus F) \ne \emptyset$. Since $V(K')$ contains a vertex of $\Gamma$, and $\Delta \setminus F$ is disjoint from $V(\Gamma)$, it follows that $U(K') \nsubseteq \Delta \setminus F$, and since $K'$ is connected, there is an edge $e$ of $K'$ with $\bar{e} \cap F \ne \emptyset$ and with $e \cap (\Delta \setminus F) \ne \emptyset$. Choose $u \in \bar{e} \cap F$, and let $u_1 \in F_1$ with $\psi(u_1) = u$. Since $u \in U(K')$, it follows that $u_1 \in bd(\Sigma_1) \subseteq U(\Gamma_1)$. Since $F_1$ is $\Gamma_1$-normal, we deduce that $u_1 \in V(\Gamma_1)$. Consequently $u \in V(\Gamma) \cup V(K)$, and since also $u \in U(K)$ and $V(\Gamma) \cap U(K) \subseteq V(K)$, it follows that $u \in V(K)$. Since $u_1$ belongs to both $F_1 \cap V(\Gamma_1)$ and $F_1 \cap bd(\Sigma_1)$, and $|F_1 \cap V(\Gamma_1)| + |F_1 \cap bd(\Sigma_1)| \le 2$, we deduce that $F_1 \cap V(\Gamma_1) = F_1 \cap bd(\Sigma_1) = \{u_1\}$; and so $F \cap (U(K') \cup bd(\Sigma)) = \{u\}$.

Let $L_1 \subseteq F_1$ be a closed line segment with $u_1$ in its interior. Since $F_1$ is an $O$-arc, there is a second line segment $L'_1 \subseteq \Sigma_1$ with the same ends as $L_1$, such that $L_1 \cup L'_1$ is an $O$-arc in $\Sigma_1$, bounding a disc $D_1 \subseteq \Sigma_1$ with $D_1 \cap F_1 = L$. Let $D = \psi(D_1), L = \psi(L_1)$, and $L' = \psi(L'_1)$. Then $D$ is a disc in $\Sigma$, with boundary $L \cup L'$, where $L$ is a line segment in $F$ with $u$ in its interior, and $D \cap F = L$. Moreover, the interior of $D$ is disjoint from $U(K)$, and in particular, disjoint from $e$. It follows that $D \cap \Delta = L$.
